\chapter{Arbitrage statistique sur résidus et composantes principales}\label{ch:s1:residual-and-principal-component-stat-arb}

On régresse les soixante derniers rendements quotidiens de chaque action sur son fonds sectoriel, on additionne les résidus et on négocie leur somme lorsqu'elle s'est nettement écartée de son niveau habituel~:
on achète lorsqu'elle se trouve à 1.25 écarts-types en dessous de sa moyenne, on vend lorsqu'elle se trouve à 1.25 écarts-types au-dessus, et on clôture près du centre. Avellaneda et Lee ont publié cette règle, avec tous ses
paramètres, et ont également rapporté un ratio de Sharpe après coûts de 1.1 entre 1997 et 2007 avec des fonds sectoriels, et de 1.44 avec des composantes principales, et ont également rapporté la
dégradation~: 0.9 entre 2003 et 2007. La méthode est publique~; les paramètres sont publiés~; ce qui détermine le résultat est de savoir si les résidus reviennent réellement à la moyenne.
Sur le marché synthétique, la même règle dégage un ratio de Sharpe net de 3.0 lorsque les actions portent une composante de retour à la moyenne, et rien lorsqu'elles n'en portent pas. La construction est
\texttt{firm.residarb}.

\section{Résidus par rapport aux facteurs ou aux composantes principales}

La stratégie part de là où le chapitre~2 s'achève~: des rendements résiduels, mais sur plusieurs semaines plutôt que sur une journée. Pour chaque action~$i$ et chaque jour, les soixante derniers
rendements quotidiens sont régressés sur les rendements des facteurs,
\[
R_{i,t} = \beta_{i,0} + \sum_j \beta_{ij} F_{j,t} + \epsilon_{i,t},
\]
et les résidus $\epsilon$ sont la part des mouvements de l'action que ses facteurs n'expliquent pas. Avellaneda et Lee ont utilisé deux types de facteurs. Dans la version «~ETF~»,
le fonds sectoriel de l'action, soit un régresseur par action. Dans la version «~PCA~», les composantes principales de l'ensemble du marché.

\begin{definition}[Portefeuille propre]\label{def:s1:residual-and-principal-component-stat-arb:eigenportfolio}
Un \emph{portefeuille propre}\index{portefeuille propre} (eigenportfolio) est le portefeuille qui investit dans chaque action l'élément correspondant d'un vecteur propre de la matrice de corrélation des
rendements des actions, divisé par la volatilité de l'action, $Q^{(j)}_i = v^{(j)}_i/\sigma_i$~; les rendements des principaux portefeuilles propres servent de facteurs à un
modèle à facteurs statistiques.
\end{definition}

La matrice de corrélation est estimée sur les 252 derniers jours~; Avellaneda et Lee ont retenu quinze \omterm{def:s1:residual-and-principal-component-stat-arb:eigenportfolio}{portefeuilles propres}, ou autant qu'il en faut pour expliquer 55~\% de la variance. Sur le
marché synthétique, dont la structure réelle compte quatorze facteurs (le marché, dix industries, trois styles), quinze \omterm{def:s1:residual-and-principal-component-stat-arb:eigenportfolio}{portefeuilles propres} expliquent 45.7~\% de la
variance, et le premier à lui seul 26.1~\%.

\section{Le résidu comme processus d'Ornstein--Uhlenbeck~: le s-score}

On additionne les soixante résidus pour obtenir une trajectoire $X_1, \dots, X_{60}$, le mouvement idiosyncratique cumulé de l'action. Si l'action a une fair value par rapport à ses
facteurs, la trajectoire devrait osciller autour de ce niveau et y revenir~: un processus d'Ornstein--Uhlenbeck (livre~4, chapitre~4), qui, échantillonné quotidiennement, est un AR(1),
\[
X_{n+1} = a + b\,X_n + \zeta_{n+1},\qquad \kappa = -252\ln b,\qquad m = \frac{a}{1-b},\qquad \sigma_{\mathrm{eq}} = \sqrt{\frac{\operatorname{Var}\zeta}{1-b^2}}.
\]
Ici $\kappa$ est la vitesse de retour à la moyenne (une année divisée par elle donne le temps de retour à la moyenne), $m$ le niveau vers lequel la trajectoire revient et $\sigma_{\mathrm{eq}}$
son écart-type autour de ce niveau.

\begin{definition}[s-score]\label{def:s1:residual-and-principal-component-stat-arb:sscore}
Le \emph{s-score}\index{s-score} d'une action est la distance de son résidu cumulé au niveau vers lequel il revient, exprimée en unités de son
écart-type d'équilibre, $s = (X - m)/\sigma_{\mathrm{eq}}$, à partir d'un ajustement AR(1) du résidu cumulé sur une fenêtre glissante. Comme les résidus de la régression ont pour somme
zéro, le dernier point vaut $X_{60} = 0$ et $s = -m/\sigma_{\mathrm{eq}}$~; Avellaneda et Lee ont centré $m$ sur l'ensemble des actions.
\end{definition}

Un \omterm{def:s1:residual-and-principal-component-stat-arb:sscore}{s-score} fortement positif signifie que l'action est chère par rapport à ses facteurs, et un \omterm{def:s1:residual-and-principal-component-stat-arb:sscore}{s-score} fortement négatif qu'elle est bon marché. La règle de négociation est une bande. On achète
(l'action, et on vend ses bêtas dans les facteurs) lorsque $s < -1.25$~; on vend lorsque $s > 1.25$~; on clôture une position longue lorsque $s > -0.50$ et une position courte lorsque $s < 0.75$. Les seuils ont été
choisis sur 2000--2004. Chaque position représente la même fraction du capital (la règle «~bang-bang~» de l'article), et chaque transaction coûte cinq points de base.

\section{Entrée, sortie et filtre de vitesse de retour à la moyenne}

\begin{definition}[Filtre de vitesse de retour à la moyenne]\label{def:s1:residual-and-principal-component-stat-arb:speed}
Un \emph{filtre de vitesse de retour à la moyenne}\index{filtre de vitesse de retour à la moyenne} ne retient le \omterm{def:s1:residual-and-principal-component-stat-arb:sscore}{s-score} d'une action que lorsque sa vitesse de retour à la moyenne ajustée est assez élevée
pour que le retour soit attendu bien avant la fin de la fenêtre d'estimation~; Avellaneda et Lee exigeaient $\kappa > 252/30 = 8.4$, un temps de retour à la moyenne inférieur à la moitié de la
fenêtre de soixante jours ($0 < b < 0.9672$).
\end{definition}

Pour tester la méthode, le marché synthétique doit contenir quelque chose qu'elle puisse trouver. \texttt{MarketConfig.ou\_share} (nouveau dans ce chapitre, désactivé par défaut afin que chaque
panel précédent reste inchangé) attribue au rendement spécifique de chaque action un niveau de retour à la moyenne qui absorbe une part de sa variance spécifique, avec des temps de retour à la moyenne
répartis selon une loi lognormale autour de vingt jours. La force de rappel attendue du niveau implanté est enregistrée, comme tout effet implanté, afin que le \omterm{def:s1:residual-and-principal-component-stat-arb:sscore}{s-score} puisse être évalué par rapport à elle~: sa
corrélation avec cette force de rappel vaut $-0.128$ pour une part de 0.3 (le signe est le bon~: un \omterm{def:s1:residual-and-principal-component-stat-arb:sscore}{s-score} élevé annonce une baisse à venir).

\begin{center}
\small
\begin{tabular}{@{}lcccccc@{}}
\toprule
part de retour à la moyenne & \multicolumn{2}{c}{0} & \multicolumn{2}{c}{0.15} & \multicolumn{2}{c}{0.3}\\
facteurs & sectoriels & PCA & sectoriels & PCA & sectoriels & PCA\\
\midrule
ratio de Sharpe avant coûts & 1.04 & 1.30 & 1.92 & 2.56 & 3.15 & 4.34\\
ratio de Sharpe après coûts & 0.07 & $-0.06$ & 0.93 & 1.22 & 2.15 & 3.00\\
rendement net annuel & 0.9~\% & $-0.5~\%$ & 11.5~\% & 9.0~\% & 26.0~\% & 22.4~\%\\
coûts annuels & 12.0~\% & 9.8~\% & 12.3~\% & 9.8~\% & 12.2~\% & 10.0~\%\\
volatilité & 12.3~\% & 7.2~\% & 12.4~\% & 7.3~\% & 12.1~\% & 7.5~\%\\
positions ouvertes & 347 & 343 & 344 & 340 & 336 & 338\\
taux de rotation quotidien & 37~\% & 39~\% & 37~\% & 39~\% & 37~\% & 40~\%\\
\bottomrule
\end{tabular}
\end{center}

Trois points ressortent. Sans composante de retour à la moyenne, la méthode couvre ses coûts et rien de plus~; le ratio de Sharpe brut d'environ 1 provient du
retournement à un jour que le marché introduit de toute façon, et cinq points de base par transaction le grignotent. Avec une telle composante, la méthode la trouve, et les composantes principales devancent les
indices sectoriels, parce qu'elles retirent aussi les facteurs de style~: la volatilité du portefeuille PCA est de 7.5~\% contre 12.1~\%. Et le portefeuille est vaste~: environ 340 positions de 1~\% du
capital chacune, une exposition brute de 3.4, proche du «~2+2~» de l'article.

Le filtre de vitesse filtre à peine. Au seuil d'Avellaneda et Lee, il retient 99.3~\% des ajustements, et il en retient 99.1~\% lorsqu'il n'y a aucun retour à la moyenne~:
un AR(1) ajusté sur soixante points d'une trajectoire ancrée à zéro en son extrémité est fortement biaisé vers un retour rapide, si bien qu'une marche aléatoire elle-même paraît rapide. Des seuils plus
stricts rejettent davantage de titres mais n'aident pas (voir la \cref{fig:s1:residual-and-principal-component-stat-arb:sens})~: exiger un temps de retour à la moyenne inférieur à quinze
jours conserve 87.9~\% des ajustements et dégage 2.03~; sous dix jours, 69.2~\% et 1.81~; sous cinq, 26.8~\% et 0.74. Sans filtre, le portefeuille sectoriel dégage 2.16 et
le portefeuille PCA dégage 3.04, contre 2.15 et 3.00 avec le filtre. Le seuil d'entrée compte peu, ce qui est une bonne nouvelle pour une règle publiée~: de 0.75 à 1.5, le ratio de Sharpe net
reste compris entre 2.20 et 2.30, et il tombe à 1.87 à 2.0, où le portefeuille compte 101 positions au lieu de 336.

\begin{omfigure}
\begin{tikzpicture}
\begin{axis}[name=a,width=6.6cm,height=5.0cm,xlabel={\small seuil d'entrée $|s|$},ylabel={\small ratio de Sharpe},tick label style={font=\footnotesize},
  xtick={0.75,1,1.25,1.5,2},ymin=0,ymax=4,grid=major,grid style={gray!25},legend pos=south west,legend style={font=\scriptsize},table/col sep=comma]
\addplot[omDef,very thick,mark=*] table[x=s_open,y=sr_gross]{figdata/strategies-1/03-residual-and-principal-component-stat-arb/sens_open.csv};
\addplot[omThm,very thick,mark=square*] table[x=s_open,y=sr_net]{figdata/strategies-1/03-residual-and-principal-component-stat-arb/sens_open.csv};
\legend{avant coûts,après coûts}
\end{axis}
\begin{axis}[at={(a.east)},anchor=west,xshift=1.3cm,width=6.6cm,height=5.0cm,xlabel={\small temps de retour à la moyenne\\maximal autorisé (jours)},xlabel style={align=center},
  ylabel={\small ratio de Sharpe net},tick label style={font=\footnotesize},xtick={5,10,15,30},xmin=2,xmax=38,ymin=0,ymax=2.6,grid=major,grid style={gray!25},
  table/col sep=comma]
\addplot[omThm,very thick,mark=square*,nodes near coords,point meta=explicit symbolic,every node near coord/.style={font=\scriptsize,anchor=north west}]
  table[x=max_tau_days,y=sr_net,meta=passed_pct]{figdata/strategies-1/03-residual-and-principal-component-stat-arb/sens_kappa.csv};
\end{axis}
\end{tikzpicture}

\omcaption{Le portefeuille de \omterm{def:s1:residual-and-principal-component-stat-arb:sscore}{s-score} sur indices sectoriels avec une part implantée de retour à la moyenne de 0.3. À gauche~: ratio de Sharpe en fonction du seuil d'entrée. À droite~: ratio de Sharpe net
en fonction du temps de retour à la moyenne maximal autorisé par le filtre de vitesse (30 jours pour Avellaneda et Lee)~; les étiquettes donnent le pourcentage d'ajustements conservés. Données~:
\texttt{s1\_residarb.run}.}\label{fig:s1:residual-and-principal-component-stat-arb:sens}
\end{omfigure}

\begin{omfigure}
\begin{tikzpicture}
\begin{axis}[width=12cm,height=4.8cm,xlabel={\small année},ylabel={\small s-score},tick label style={font=\footnotesize},xmin=5,xmax=6,ymin=-3,ymax=3,
  ytick={-1.25,0,1.25},grid=major,grid style={gray!25},table/col sep=comma]
\addplot[omDef,thick] table[x=year,y=s]{figdata/strategies-1/03-residual-and-principal-component-stat-arb/path.csv};
\addplot[omThm,thick,const plot] table[x=year,y expr=\thisrow{pos}*2.6]{figdata/strategies-1/03-residual-and-principal-component-stat-arb/path.csv};
\addplot[black!50,dashed] coordinates {(5,1.25) (6,1.25)};
\addplot[black!50,dashed] coordinates {(5,-1.25) (6,-1.25)};
\end{axis}
\end{tikzpicture}

\omcaption{Le \omterm{def:s1:residual-and-principal-component-stat-arb:sscore}{s-score} d'une action synthétique (par rapport à son indice sectoriel, non centré) sur une année, les seuils d'entrée étant tracés en pointillés~; la ligne en escalier représente la position
(montée en position longue, descente en position courte, tracée à $\pm 2.6$). Données~: \texttt{s1\_residarb.path}.}\label{fig:s1:residual-and-principal-component-stat-arb:path}
\end{omfigure}

\section{Portefeuilles propres et leur stabilité}

Un facteur qui change chaque mois est un facteur que le portefeuille ne peut pas couvrir à bon compte. Réestimé tous les vingt et un jours sur une année glissante, le premier vecteur propre est
presque identique à chaque fois (le cosinus en valeur absolue entre estimations successives vaut en moyenne 0.999~: c'est le marché). Le sous-espace de dimension quinze bouge davantage~:
son recouvrement moyen avec celui du mois précédent (le cosinus carré moyen des angles principaux) vaut 0.89, et le cosinus de la direction la moins stable vaut en moyenne 0.63
(\cref{fig:s1:residual-and-principal-component-stat-arb:stab}). Les composantes suivantes mélangent les industries et les styles de façon différente d'une estimation à l'autre, et
une couverture dans ces composantes est reconstruite chaque mois. La version fondée sur l'indice sectoriel paie sa stabilité par un ajustement moins bon~; la version en ACP paie son ajustement par des couvertures qui se
renouvellent, ce que le modèle de coûts de ce chapitre facture à cinq points de base par unité de chaque \omterm{def:s1:residual-and-principal-component-stat-arb:eigenportfolio}{portefeuille propre}, une simplification généreuse pour un panier de
mille titres.

\begin{omfigure}
\begin{tikzpicture}
\begin{axis}[width=10cm,height=5.2cm,xlabel={\small année},ylabel={\small recouvrement avec le précédent},tick label style={font=\footnotesize},xmin=2,xmax=10,
  ymin=0,ymax=1,grid=major,grid style={gray!25},legend style={font=\scriptsize,at={(0.5,1.03)},anchor=south,legend columns=2},
  table/col sep=comma]
\addplot[omDef,very thick] table[x=year,y=subspace]{figdata/strategies-1/03-residual-and-principal-component-stat-arb/stability.csv};
\addplot[omThm,thick] table[x=year,y=min_cos]{figdata/strategies-1/03-residual-and-principal-component-stat-arb/stability.csv};
\legend{sous-espace de dimension 15,direction la moins stable}
\end{axis}
\end{tikzpicture}

\omcaption{Stabilité des quinze premiers \omterm{def:s1:residual-and-principal-component-stat-arb:eigenportfolio}{portefeuilles propres}, réestimés tous les 21 jours sur les 252 derniers~: le recouvrement avec l'estimation précédente du sous-espace
qu'ils engendrent (cosinus carré moyen des angles principaux) et le cosinus du plus grand angle principal. Données~: \texttt{s1\_residarb.stability}.}
\label{fig:s1:residual-and-principal-component-stat-arb:stab}
\end{omfigure}

\section{Où sont passés les rendements}

Les chiffres mêmes d'Avellaneda et Lee racontent le déclin de la stratégie~: le ratio de Sharpe net de la version en ACP s'établit en moyenne à 1.44 sur 1997--2007, nettement
plus élevé avant 2003 et de 0.9 de 2003 à 2007~; la version ETF affiche 1.1 sur l'ensemble de la période, avec une dégradation comparable après 2002. Deux autres constats du même
article indiquent ce qui a continué à fonctionner. En prenant le temps de transaction pour référence, en divisant le rendement de chaque jour par le volume de ce jour rapporté à sa moyenne sur dix jours, on porte la version ETF
à 1.51 sur 2003--2007. Le comportement des stratégies à l'été 2007 correspond au récit de Khandani et Lo sur des portefeuilles de ce type
débouclés simultanément (livre~7, chapitre~28).

Le marché synthétique ne peut pas dire pourquoi les résidus réels revenaient moins à la moyenne après 2002, seulement ce dont la méthode a besoin~: une composante à retour à la moyenne assez grande pour
absorber cinq points de base par transaction sur un tiers du portefeuille chaque jour. Avec une part de 0.15 de la variance spécifique, le portefeuille sectoriel affiche un ratio net de 0.93~; à 0.3, 2.15~; à zéro, rien.
Son propre effet de volume joue en sens inverse~: le volume du marché synthétique croît avec l'ampleur du choc spécifique du jour, si bien que le temps de transaction réduit précisément les
mouvements qui s'inversent, et le portefeuille sectoriel ajusté du volume affiche un ratio net de 1.78 contre 2.15. Que le volume soit utile dépend de ce qu'il mesure sur un marché donné~; l'article
a constaté que le volume avait aidé aux États-Unis sur 2003--2007.

\section{Fiches de stratégie}

\begin{strategyfile}{Stratégie de s-score sur résidus d'ETF}\label{strat:s1:residual-and-principal-component-stat-arb:etf}
\sfield{Qui paie l'arbitragiste, et pourquoi} Des flux qui écartent une action de son secteur sans nouvelle information~: transactions sur indices et fonds, vendeurs pressés~; le portefeuille les absorbe et
attend que l'action rejoigne son secteur.
\sfield{Instruments et plateformes} Actions liquides et leurs fonds sectoriels, en positions longues et courtes~; les couvertures passent par les fonds.
\sfield{Signal} Le \omterm{def:s1:residual-and-principal-component-stat-arb:sscore}{s-score} du résidu cumulé sur 60 jours par rapport au fonds sectoriel~; ouverture à $|s| > 1.25$, clôture à $-0.50$ (positions longues) et $0.75$ (positions courtes).
\sfield{Dimensionnement et exécution} Positions égales d'une fraction fixe du capital~; chacune couverte par son bêta dans le fonds~; couvertures compensées sur l'ensemble du portefeuille.
\sfield{Coûts} Cinq points de base par transaction dans l'article~; un tiers du portefeuille négocié chaque jour ici.
\sfield{Comment elle s'éteint} Les résidus cessent de revenir à la moyenne (le bilan publié après 2002)~; encombrement (crowding) et débouclages simultanés (août 2007)~; coûts supérieurs à l'edge.
\sfield{Horizon, capacité, infrastructure} De quelques jours à quelques semaines~; des centaines de positions~; une régression et un ajustement quotidiens par action, couverture par les fonds sectoriels.
\sfield{Backtest honnête} Univers et appartenance sectorielle à chaque date~; coûts sur les deux jambes~; seuils choisis sur une période antérieure à celle qui est rapportée.
\sfield{Sources} Avellaneda et Lee (2010)~: ratio de Sharpe de 1.1 après coûts sur 1997--2007, en baisse après 2002~; simulation de ce chapitre.
\end{strategyfile}

\begin{strategyfile}{Stratégie de résidus en composantes principales}\label{strat:s1:residual-and-principal-component-stat-arb:pca}
\sfield{Qui paie l'arbitragiste, et pourquoi} Comme la version ETF~; retirer davantage de structure commune laisse un résidu plus propre.
\sfield{Instruments et plateformes} Actions liquides~; couvertures dans les \omterm{def:s1:residual-and-principal-component-stat-arb:eigenportfolio}{portefeuilles propres}, qui sont des paniers des mêmes actions.
\sfield{Signal} Le \omterm{def:s1:residual-and-principal-component-stat-arb:sscore}{s-score} par rapport à 15 \omterm{def:s1:residual-and-principal-component-stat-arb:eigenportfolio}{portefeuilles propres} issus de la matrice de corrélation de l'année écoulée (ou assez pour expliquer 55~\% de la variance).
\sfield{Dimensionnement et exécution} Comme la version ETF~; couvertures compensées dans le portefeuille d'actions, puisque les \omterm{def:s1:residual-and-principal-component-stat-arb:eigenportfolio}{portefeuilles propres} sont des actions.
\sfield{Coûts} Transactions sur actions, plus la dérive des \omterm{def:s1:residual-and-principal-component-stat-arb:eigenportfolio}{portefeuilles propres} entre deux réestimations.
\sfield{Comment elle s'éteint} Comme la version ETF~; \omterm{def:s1:residual-and-principal-component-stat-arb:eigenportfolio}{portefeuilles propres} qui se renouvellent plus vite que le portefeuille ne peut les couvrir.
\sfield{Horizon, capacité, infrastructure} De quelques jours à quelques semaines~; une estimation de la matrice de corrélation et de ses vecteurs propres, rafraîchie régulièrement.
\sfield{Backtest honnête} Vecteurs propres estimés uniquement sur les données antérieures à chaque transaction~; nombre de composantes fixé à l'avance.
\sfield{Sources} Avellaneda et Lee (2010)~: ratio de Sharpe de 1.44 après coûts sur 1997--2007, 0.9 sur 2003--2007.
\end{strategyfile}

\begin{strategyfile}{Portefeuille de résidus neutre par industrie}\label{strat:s1:residual-and-principal-component-stat-arb:industry}
\sfield{Qui paie l'arbitragiste, et pourquoi} Comme précédemment, mais les \omterm{def:s1:residual-and-principal-component-stat-arb:sscore}{s-scores} sont transformés en portefeuille plutôt qu'en ensemble de transactions.
\sfield{Instruments et plateformes} Actions liquides uniquement~; aucune couverture par fonds.
\sfield{Signal} Opposé du \omterm{def:s1:residual-and-principal-component-stat-arb:sscore}{s-score}, pris comme prévision en unités de rendement (chapitre~2).
\sfield{Dimensionnement et exécution} Un optimiseur neutre par rapport aux industries et aux styles, coûts intégrés (chapitre~2), au lieu de positions égales en tout ou rien.
\sfield{Coûts} Inférieurs à la règle tout ou rien~: l'optimiseur n'échange que là où la prévision est rentable.
\sfield{Comment elle s'éteint} Comme précédemment.
\sfield{Horizon, capacité, infrastructure} De quelques jours à quelques semaines~; un modèle de risque et un optimiseur.
\sfield{Backtest honnête} Les mêmes ajustements et seuils que la version tout ou rien, de sorte que la comparaison isole la construction du portefeuille.
\sfield{Sources} Avellaneda et Lee (2010) notent que la règle du tout ou rien a surpassé les ajustements continus dans leurs tests, qu'ils attribuent à une
erreur de spécification du modèle~; la comparaison est à mener par le lecteur (tutoriel).
\end{strategyfile}

\begin{strategyfile}{Stratégie de résidus conditionnée par le volume}\label{strat:s1:residual-and-principal-component-stat-arb:volume}
\sfield{Qui paie l'arbitragiste, et pourquoi} Comme la version ETF, en pondérant les mouvements selon le volume de transactions qu'ils ont provoqué.
\sfield{Instruments et plateformes} Comme la version ETF.
\sfield{Signal} Le \omterm{def:s1:residual-and-principal-component-stat-arb:sscore}{s-score} des résidus de rendements exprimés en temps de transaction~: le rendement de chaque jour multiplié par le rapport entre son volume moyen sur dix jours et son propre volume.
\sfield{Dimensionnement et exécution} Comme la version ETF.
\sfield{Coûts} Comme la version ETF.
\sfield{Comment elle s'éteint} Quand le volume cesse de distinguer les mouvements de liquidité des mouvements d'actualité~; sur le marché synthétique, où le volume augmente avec chaque grand choc,
elle y nuit.
\sfield{Horizon, capacité, infrastructure} De quelques jours à quelques semaines~; volume quotidien par action.
\sfield{Backtest honnête} Seul le volume connu à la clôture est utilisé~; l'ajustement est comparé à la version simple sur la même période.
\sfield{Sources} Avellaneda et Lee (2010)~: 1.51 sur 2003--2007 pour la version ETF avec information sur le volume.
\end{strategyfile}

\section{Tutoriel~: la méthode publique}\label{tut:s1:residual-and-principal-component-stat-arb:tut}

\begin{tutorial}
\textbf{Objectif.} Faire tourner la règle publiée avec des indices sectoriels et avec des composantes principales, avec et sans composante de retour à la moyenne, et mesurer l'effet de ses
paramètres. \textbf{État final~:} le tableau et la \cref{fig:s1:residual-and-principal-component-stat-arb:sens}.
\begin{tutsteps}
\item \textbf{L'ajustement et la règle}~: l'AR(1) du résidu cumulé, le \omterm{def:s1:residual-and-principal-component-stat-arb:sscore}{s-score} et la bande.
\omcode{code/firm/residarb/firm_residarb.py}{51}{80}{Ajustement du processus d'Ornstein--Uhlenbeck par AR(1), s-score et règle de négociation.}\label{lst:s1:residual-and-principal-component-stat-arb:fit}
\item \textbf{Chaque jour}~: régresser, ajuster, calculer le score, négocier, couvrir.
\omcode{code/strategies-1/03-residual-and-principal-component-stat-arb/python/s1_residarb.py}{89}{108}{Étape quotidienne du portefeuille de stat arb au s-score.}\label{lst:s1:residual-and-principal-component-stat-arb:loop}
\item \textbf{Exécuter} \texttt{run(share, method, s\_open, min\_kappa, volume)} sur la grille, \texttt{stability()} et \texttt{fig\_residarb.py}.
\end{tutsteps}
\textbf{Pour aller plus loin.} Remplacer la règle bang-bang par l'optimiseur sensible aux coûts du chapitre~2 (le fichier de la stratégie neutre en secteurs)~; estimer la régression
sur 90 jours et le processus sur 60, comme le suggère l'article~; raccourcir les temps de retour à la moyenne et observer à partir de quand le filtre de vitesse commence à peser.
\end{tutorial}

\section{Construction~: stat arb sur résidus}\label{bld:s1:residual-and-principal-component-stat-arb:residarb}

\begin{build}
\bfield{Rôle} Résidus, ajustements OU, \omterm{def:s1:residual-and-principal-component-stat-arb:sscore}{s-scores} et règle de négociation de la stat arb sur résidus, pour un ensemble de facteurs quelconque.
\bfield{Interface} \texttt{regress(R, F, per\_name)}, \texttt{eigenportfolios(R, k)}, \texttt{ou\_fit(resid)}, \texttt{s\_score(fit, center)},
\texttt{step(pos, s, ok, s\_open, s\_close\_long, s\_close\_short)}, \texttt{speed\_ok(kappa, min\_kappa)}, \texttt{volume\_adjust(R, V, avg)}.
\bfield{Règles} Seules les données antérieures à chaque transaction sont utilisées~; les ajustements dont $b \notin (0, 1)$ sont rejetés~; les paramètres sont fixés à l'avance et publiés avec les résultats.
\bfield{Tests d'acceptation} \texttt{code/firm/residarb/tests/}~: les régressions retrouvent des bêtas connus~; l'AR(1) retrouve un coefficient connu,
à son biais de petit échantillon près~; le \omterm{def:s1:residual-and-principal-component-stat-arb:sscore}{s-score} et la bande vérifiés à la main~; les \omterm{def:s1:residual-and-principal-component-stat-arb:eigenportfolio}{portefeuilles propres} d'un marché à un facteur.
\bfield{Pour aller plus loin} Un terme de dérive dans le \omterm{def:s1:residual-and-principal-component-stat-arb:sscore}{s-score} (le «~\omterm{def:s1:residual-and-principal-component-stat-arb:sscore}{s-score} modifié~» de l'article)~; des estimations de l'AR(1) corrigées du biais~; des couvertures dans de vrais ETF avec leurs propres coûts.
\end{build}

\begin{omsources}
\item M.~Avellaneda et J.-H.~Lee, «~Statistical arbitrage in the US equities market~», \emph{Quantitative Finance} 10(7), 2010 (version de travail de
juin 2009 sur la page du premier auteur).
\item A.~E.~Khandani et A.~W.~Lo, «~What happened to the quants in August 2007~?~», \emph{Journal of Financial Markets} 14(1), 2011.
\item A.~Pole, \emph{Statistical Arbitrage}, Wiley.
\end{omsources}

\section{Exercices}

\begin{exercise}[$\star$]\label{exo:s1:residual-and-principal-component-stat-arb:1}
Un coefficient AR(1) du résidu cumulé vaut $b = 0.9$ par jour. Calculer $\kappa$ et le temps de retour à la moyenne en jours.
\end{exercise}

\begin{exercise}[$\star$]\label{exo:s1:residual-and-principal-component-stat-arb:2}
Avec $b = 0.9$, $a = 0.002$ et $\operatorname{Var}\zeta = 10^{-4}$, calculer $m$, $\sigma_{\mathrm{eq}}$ et le \omterm{def:s1:residual-and-principal-component-stat-arb:sscore}{s-score} (non centré). La règle déclenche-t-elle une transaction~?
\end{exercise}

\begin{exercise}[$\star$]\label{exo:s1:residual-and-principal-component-stat-arb:3}
Montrer que $\kappa > 252/30$ est la même condition que $b < 0.9672$.
\end{exercise}

\begin{exercise}[$\star\star$]\label{exo:s1:residual-and-principal-component-stat-arb:4}
Le portefeuille sectoriel fait tourner 37~\% du capital par jour. Que coûtent cinq points de base par unité négociée sur un an, pour les seules actions, et pourquoi le
chiffre de coût du chapitre est-il plus élevé~?
\end{exercise}

\begin{exercise}[$\star\star$]\label{exo:s1:residual-and-principal-component-stat-arb:5}
Pourquoi le filtre de vitesse laisse-t-il passer 99~\% des ajustements, même sur un marché sans retour à la moyenne~?
\end{exercise}

\begin{exercise}[$\star\star$]\label{exo:s1:residual-and-principal-component-stat-arb:6}
Pourquoi le portefeuille en composantes principales a-t-il une volatilité plus faible que le portefeuille sectoriel, à nombre de positions égal~?
\end{exercise}

\begin{exercise}[$\star\star\star$]\label{exo:s1:residual-and-principal-component-stat-arb:7}
\emph{Programmation.} Exécuter \texttt{run(0.15, 'pca')} et \texttt{run(0.0, 'pca')}. Donner les ratios de Sharpe nets et expliquer ce que mesure le second.
\end{exercise}

\begin{exercise}[$\star\star\star$]\label{exo:s1:residual-and-principal-component-stat-arb:8}
\emph{Trouver la faille.} «~La méthode est publiée avec ses paramètres et un ratio de Sharpe de 1.44 après coûts~; nous allons la faire tourner telle quelle.~»
\end{exercise}

\section{Problème~: la méthode publique}

\begin{problem}[{Problème du week-end --- une stratégie publiée, relancée}]\label{pb:s1:residual-and-principal-component-stat-arb:1}
La règle d'Avellaneda et Lee sur \texttt{firm.synthmkt}, avec et sans composante de retour à la moyenne injectée.

\medskip\noindent\textbf{Partie I --- La méthode.}
\begin{enumerate}
\item Décrire les deux façons de calculer les résidus.
\item Définir un \omterm{def:s1:residual-and-principal-component-stat-arb:eigenportfolio}{portefeuille propre} et indiquer combien en retient l'article.
\item Ajuster le processus OU au résidu cumulé~: quelles sont les valeurs de $\kappa$, $m$ et $\sigma_{\mathrm{eq}}$~?
\item Définir le \omterm{def:s1:residual-and-principal-component-stat-arb:sscore}{s-score} et donner les seuils de la règle de négociation.
\end{enumerate}

\medskip\noindent\textbf{Partie II --- Le test synthétique.}
\begin{enumerate}[resume]
\item Que \texttt{ou\_share} injecte-t-il, et comment le \omterm{def:s1:residual-and-principal-component-stat-arb:sscore}{s-score} est-il évalué par rapport à cette injection~?
\item Donner les ratios de Sharpe nets des deux versions pour des parts de 0, 0.15 et 0.3.
\item Pourquoi la version en composantes principales fait-elle mieux ici~?
\item Que rapporte la méthode sans retour à la moyenne, et d'où vient son rendement brut~?
\end{enumerate}

\medskip\noindent\textbf{Partie III --- Les paramètres.}
\begin{enumerate}[resume]
\item Quelle part des ajustements passe le filtre de vitesse de l'article, avec et sans retour à la moyenne, et pourquoi~?
\item Que font des filtres plus stricts~?
\item Quelle est la sensibilité du résultat au seuil d'entrée~?
\item Que fait ici l'ajustement de volume, et pourquoi~?
\end{enumerate}

\medskip\noindent\textbf{Partie IV --- Le verdict.}
\begin{enumerate}[resume]
\item Énoncer le \emph{résultat nommé}~: le ratio de Sharpe de la stratégie au \omterm{def:s1:residual-and-principal-component-stat-arb:sscore}{s-score} avec et sans filtre de vitesse, et sa sensibilité au seuil d'entrée.
\item Quelle est la stabilité des \omterm{def:s1:residual-and-principal-component-stat-arb:eigenportfolio}{portefeuilles propres}~?
\item Que l'article a-t-il rapporté après 2002~?
\item Qu'apporte août 2007~?
\item De quoi la méthode a-t-elle besoin d'un marché pour fonctionner~?
\item Comment saurait-on, sur des données réelles, si les résidus reviennent à la moyenne~?
\item Quel paramètre inquiéterait le plus~?
\item En une phrase~: qu'est-ce que la stat arb sur résidus~?
\end{enumerate}
\end{problem}

\section{Questions d'entretien}

\begin{interviewq}[$\star$ \iqroles{chercheur, trader}]\label{iq:s1:residual-and-principal-component-stat-arb:1}
Qu'est-ce qu'un \omterm{def:s1:residual-and-principal-component-stat-arb:sscore}{s-score}, et comment le négocieriez-vous~?
\end{interviewq}

\begin{interviewq}[$\star\star$ \iqroles{chercheur}]\label{iq:s1:residual-and-principal-component-stat-arb:2}
Facteurs en composantes principales ou ETF sectoriels pour la stat arb sur résidus~: quels sont les compromis~?
\end{interviewq}

\begin{interviewq}[$\star\star$ \iqroles{chercheur}]\label{iq:s1:residual-and-principal-component-stat-arb:3}
Vous ajustez un AR(1) sur 60 observations et obtenez $b = 0.85$. Quelle confiance accordez-vous au temps de retour à la moyenne~?
\end{interviewq}

\begin{interviewq}[$\star\star$ \iqroles{risque}]\label{iq:s1:residual-and-principal-component-stat-arb:4}
Un portefeuille de stat-arb sur résidus perd 5~\% en trois jours sans nouvelle. Que vérifiez-vous~?
\end{interviewq}

\begin{interviewq}[$\star\star$ \iqroles{développeur}]\label{iq:s1:residual-and-principal-component-stat-arb:5}
Calculez les \omterm{def:s1:residual-and-principal-component-stat-arb:sscore}{s-scores} de 3\,000 actions chaque jour en moins d'une seconde~: comment faire~?
\end{interviewq}

\begin{interviewq}[$\star\star\star$ \iqroles{chercheur}]\label{iq:s1:residual-and-principal-component-stat-arb:6}
Dérivez l'écart-type d'équilibre d'un processus OU échantillonné chaque jour comme un AR(1), et montrez pourquoi $X_{60} = 0$ dans le \omterm{def:s1:residual-and-principal-component-stat-arb:sscore}{s-score} de l'article.
\end{interviewq}
